\documentclass{ieeeaccess}
\usepackage{cite}
\usepackage{amsmath,amssymb,amsfonts}
\usepackage{algorithmic}
\usepackage{graphicx}
\usepackage{textcomp}
\usepackage{booktabs}
\usepackage{tabularx}
\usepackage{array}
\usepackage{colortbl}
\definecolor{execbg}{rgb}{0.90,0.95,1.00}
\definecolor{layerbg}{rgb}{0.94,0.94,0.94}
\definecolor{crossbg}{rgb}{1.00,0.93,0.86}
\setlength{\emergencystretch}{2em}
\hbadness=10000
\vbadness=10000
\hfuzz=510pt
\AtBeginDocument{%
\setlength{\emergencystretch}{2em}%
\hbadness=10000%
\vbadness=10000%
\hfuzz=510pt%
}
\newdimen\xfigwd
\xfigwd=\columnwidth
\providecommand{\euro}{\texteuro}
\def\BibTeX{{\rm B\kern-.05em{\sc i\kern-.025em b}\kern-.08em
    T\kern-.1667em\lower.7ex\hbox{E}\kern-.125emX}}
\begin{document}
\history{Date of publication xxxx 00, 0000, date of current version xxxx 00, 0000.}
\doi{10.1109/ACCESS.2017.DOI}

\title{Cybersecurity and Cyber-Resilience for the European Cloud--Edge--IoT
Computing Continuum: A Strategic Capability Roadmap}
\author{\uppercase{First A. Author}\authorrefmark{1}, \IEEEmembership{Member, IEEE},
\uppercase{Second B. Author}\authorrefmark{2}, and
\uppercase{Third C. Author}\authorrefmark{3}}
\address[1]{Affiliation 1 (e-mail: author1@domain)}
\address[2]{Affiliation 2 (e-mail: author2@domain)}
\address[3]{Affiliation 3 (e-mail: author3@domain)}
\tfootnote{This work was supported in part by the NexusForum coordination
action and the EuCloudEdgeIoT initiative under Horizon Europe.}

\markboth
{Author \headeretal: Cybersecurity for the European Cloud--Edge--IoT Continuum}
{Author \headeretal: Cybersecurity for the European Cloud--Edge--IoT Continuum}

\corresp{Corresponding author: First A. Author (e-mail: author1@domain).}

\begin{abstract}
The distribution of computation across cloud, edge, and IoT tiers, combined
with the simultaneous operationalisation of NIS2, the Cyber Resilience Act,
the AI Act, and DORA, has made cybersecurity, resilience, and digital
sovereignty an integrated strategic problem for Europe. Perimeter-oriented
security models do not scale to heterogeneous, physically exposed, and
multi-jurisdictional continuum infrastructures. This paper addresses that
problem by proposing a three-stage methodology for identifying and
prioritising cybersecurity capabilities: architectural decomposition into
six Security Layers and four Transversal Cross-Layers, seven-dimension
capability profiling, and priority clustering with phased transition
matrices. Applying the methodology, the paper identifies ten strategic
capabilities spanning hardware roots of trust, post-quantum cryptographic
agility, lightweight edge security, federated interoperability, supply chain
verification, cognitive security operations, pan-European network
federation, confidential computing, resilient infrastructure, and proactive
threat management. These are consolidated into three executive priority
clusters, Sovereign Trusted Infrastructure, Cognitive Security and
Autonomous Cyber Defence, and Federated Operational Sovereignty, analysed
over a 2026--2040 transition horizon. The contribution is both methodological,
providing a reproducible template for cybersecurity roadmap construction,
and strategic, framing autonomous, federated, and certified
Security-by-Design as the operational target for European digital
infrastructure.
\end{abstract}

\begin{keywords}
Certification, Cloud--Edge--IoT continuum, confidential computing,
cyber-resilience, digital sovereignty, European Union, post-quantum
cryptography, security roadmap, Zero-Trust architecture.
\end{keywords}

\titlepgskip=-15pt
\maketitle

\section{Introduction}
\label{sec:intro}

The shift from centralised cloud to a distributed Cloud--Edge--IoT continuum
relocates computation, storage, and inference closer to devices, users, and
physical processes \cite{dustdar2023,bittencourt2018,ecroadmap2024}. The
architectural benefit is lower latency and higher operational autonomy; the
security consequence is a sharp increase in attack surface,
heterogeneity, and the number of administrative and trust domains that a
defender must reason about simultaneously. Perimeter-oriented and siloed
security models, designed for consolidated data-centre environments, are
not structurally adequate for infrastructures that are geographically
dispersed, resource-constrained, and physically exposed
\cite{rose2020zta,chandramouli2023zta}.

Europe faces this transition under four simultaneous regulatory
activations. NIS2 \cite{eunis2}, the Cyber Resilience Act (CRA)
\cite{eucra2024}, the AI Act \cite{euaiact}, and DORA \cite{eudora}
together raise the baseline of mandatory security, resilience,
accountability, and supply chain obligations across essential entities,
ICT products, high-risk AI systems, and financial sector third-party
providers. The EU cybersecurity certification framework
\cite{eucsa,commreg2024eucc} provides the assurance substrate these
regulations presuppose. Operationalisation of this combined regime is
uneven: national transposition lags the legal timetable in several large
Member States \cite{ec2025nis2status}, EUCS remains unadopted
\cite{ec2025cloudsov}, and conformity assessment infrastructure for
AI-specific security obligations is not yet in place \cite{ep2026aiact}.

The central argument of this paper is that Europe's normative leadership
in trust, privacy, governance, and certification
\cite{madiega2020,farrand2024} translates into strategic advantage only
when it is embedded in deployable capabilities. We call the distance
between normative leadership and deployable capacity the
\emph{ambition--implementation gap}. Closing this gap requires a
transition from reactive, siloed security towards autonomous, federated,
and certified Security-by-Design: infrastructures that sustain
trustworthy operation under adverse conditions (autonomous), security
mechanisms that span providers and jurisdictions (federated), and
assurance embedded through the system lifecycle (certified
Security-by-Design) \cite{enisa2019iotsdlc,cavoukian2009,rose2020zta}.

The paper makes four contributions. First, it frames continuum
cybersecurity as an interlocking problem across six Security Layers and
four Transversal Cross-Layers (Section~\ref{sec:model}). Second, it
defines a three-stage methodology linking architectural decomposition to
European strategic concerns through a Sovereignty Level construct
(Section~\ref{sec:method}). Third, it identifies ten strategic
capabilities, consolidated into three priority clusters analysed over a
2026--2040 transition horizon (Sections~\ref{sec:focus}--\ref{sec:clusters}).
Fourth, it derives implementation-oriented recommendations connecting
research, policy, industrial capacity, and sovereignty
(Section~\ref{sec:recs}).

Section~\ref{sec:background} provides the strategic and regulatory
context. Section~\ref{sec:model} defines the system model.
Section~\ref{sec:method} presents the methodology.
Section~\ref{sec:focus} identifies focus areas.
Section~\ref{sec:clusters} presents the three priority clusters and
their transition matrices. Section~\ref{sec:recs} states the
recommendations. Section~\ref{sec:disc} discusses validation and
limitations. Section~\ref{sec:concl} concludes.

\section{Background and Driving Factors}
\label{sec:background}

The continuum is characterised by decentralisation and physical
exposure. Decentralisation complicates visibility, orchestration,
patching, identity, and incident response across multiple administrative
domains and increases the probability of inconsistent security postures
across deployments \cite{roman2018,ranaweera2021}. Physical exposure
removes the environmental and access controls that hyperscale data
centres rely on: edge nodes in factories, substations, transport hubs,
and telecom sites are exposed to tampering, fault injection, side-channel
attacks, and supply chain compromise \cite{alwarafy2021}. Hardware roots
of trust, secure boot, remote attestation, and tamper-aware design are
baseline requirements rather than optional enhancements
\cite{tpm2spec,sailer2004}.

The European regulatory activation is historically dense. NIS2
\cite{eunis2}, the CRA \cite{eucra2024}, the AI Act \cite{euaiact}, DORA
\cite{eudora}, the Data Act \cite{eudataact}, and GDPR \cite{eugdpr}
together impose obligations that must be translated into architectures,
assurance mechanisms, and lifecycle controls operating across cloud,
edge, and IoT. The simultaneity burdens SMEs and smaller Member States
disproportionately.

Sovereignty is the third vector. Dependence on non-EU providers for
secure silicon, cloud platforms, and enterprise security tooling creates
geopolitical and operational risk. Extra-territorial data access under
foreign jurisdiction has been confirmed as a practical constraint for EU
users of hyperscaler services \cite{microsoftsenat2025}.

Europe's position in this landscape is asymmetric.
A world-leading regulatory foundation and a rigorous certification
framework \cite{commreg2024eucc,eucsa} coexist with structural industrial
gaps. EU-headquartered vendors (Infineon, NXP, STMicroelectronics) design
the dominant secure element products and achieved EUCC ``High''
certification for multiple product families in 2025
\cite{enisaeucc2025}, but fabrication depends on non-EU foundries and
designs remain proprietary, preventing sovereign auditability at the
silicon level. Intel, AMD, and ARM control the operational Trusted
Execution Environment (TEE) market. No EU-origin RISC-V processor has
received EUCC or Common Criteria certification as of April 2026. The
enterprise SIEM market is dominated by US platforms, with no
EU-headquartered vendor in the Gartner Leaders quadrant as of late 2025
\cite{gartner2025siem}. Hyperscalers hold approximately 70\% of the EU
cloud market; EU-headquartered providers collectively hold about 15\%,
stable since 2022 but below the 29\% share recorded in 2017
\cite{synergy2025}. For post-quantum migration the NIS Cooperation Group
roadmap \cite{niscg2025pqc} now provides a binding framework, but
production-scale post-quantum cryptographic (PQC) hardware tailored for
edge and IoT is not yet available in Europe.

\section{Methodology}
\label{sec:method}

The methodology transforms a heterogeneous strategic landscape into a structured roadmap for research, innovation, policy, deployment, and ecosystem development. It proceeds through three staged phases (architectural decomposition; capability formulation and profiling; priority clustering with phased progression), followed by iterative validation.

Two commitments shape the methodology. First, it supports strategic roadmap construction, not deterministic forecasting. Progress in complex socio-technical ecosystems is non-linear and depends on interacting technical, regulatory, operational, and market developments whose timing cannot be predicted precisely. Second, the methodology connects technical cybersecurity capabilities with broader European strategic concerns: resilience, interoperability, trust, and \emph{Sovereignty Level}, understood as the degree to which critical capabilities, control points, and dependencies remain under credible European influence, governance, or operational control. Sovereignty Level is an interpretive strategic lens rather than a rigid scoring model.

\subsection{Stage 1: Core Dimensions and Focus Area Identification}

The first stage decomposes the Cloud--Edge--IoT continuum into architectural vertical and cross-cutting layers, treated as analytically distinct but operationally interdependent. A vulnerability or governance gap at one layer typically propagates to adjacent layers, and effective security requires coordinated action across the full vertical and horizontal extent of the architecture.

Stage 1 answers three questions: Where are the major security and resilience problem domains within the continuum? At what architectural layers do these problems emerge? Which focus areas are sufficiently strategic, cross-cutting, or underdeveloped to justify dedicated roadmap attention?

Focus areas are derived from four sources: the architectural characteristics of distributed systems (heterogeneity across HPC, cloud, edge, and device tiers creates structurally distinct security requirements); the recurring cybersecurity challenges observed across these tiers (hardware integrity, software supply chain verification, workload isolation, identity federation, data confidentiality, network perimeter control); cross-cutting strategic concerns (sovereignty, trust, interoperability, resilience, assurance, regulatory alignment); and the need to connect technical architecture with operational and governance realities.

The output of Stage 1 is a structured map of the problem space, organised into six Security Layers and four Transversal Cross-Layers (Section~\ref{sec:sysmodel}). Sovereignty-sensitive focus areas, where dependencies on external providers or assurance mechanisms are structurally significant, are identified explicitly so that subsequent stages can track their evolution.

\subsection{Stage 2: Capability Identification and Profiling}

A capability is a strategic ability that the European ecosystem should be able to design, deploy, govern, standardise, assure, coordinate, or scale in order to achieve trustworthy, resilient, and sovereign operation of the continuum. Capabilities may encompass technical controls, governance mechanisms, interoperability instruments, assurance functions, ecosystem coordination capacities, and deployment enablers.

Capabilities are derived by translating each focus area identified in Stage 1 into actionable strategic needs, guided by a consistent question: what must Europe be able to do, technically, institutionally, and industrially, to address this focus area at the scale and assurance level that the regulatory and threat environment demands? The ten capabilities identified span hardware roots of trust, post-quantum cryptographic agility, lightweight edge security, federated interoperability, supply chain verification, cognitive security operations, pan-European network federation, confidential computing, resilient infrastructure, and proactive threat management.

Each capability is assessed through a seven-dimension template applied identically across all priorities: (D1) Baseline and Target Objective; (D2) Primary Drivers, distinguishing threat from regulatory drivers; (D3) High-Impact Activities, with three to five concrete research and innovation actions tagged by target phase; (D4) Catalysts and Enablers, covering funding alignment, regulatory and standardisation mechanisms, and ecosystem enablers; (D5) Disruptors and Systemic Risks, covering technical cliffs, supply chain risks, and geopolitical risks; (D6) Failsafe Mechanisms, specifying engineering fallbacks; and (D7) Transversal Meta-Layer, connecting the capability to sovereignty, trust and ethics, and sustainable security axes. The template enables structured cross-priority comparison, essential both for funding prioritisation and for identifying shared dependencies.

\subsection{Stage 3: Priority Clustering and Transition Matrix}

Capabilities that share architectural proximity, common enabling conditions, or convergent strategic outcomes benefit from coordinated treatment. Stage 3 groups the ten capabilities into three Priority Clusters:
(i) \emph{Sovereign Trusted Infrastructure}, the hardware, cryptographic, and resilience foundations that anchor trust across the continuum;
(ii) \emph{Cognitive Security and Autonomous Cyber Defence}, the AI-driven operational layer that scales defence beyond human-paced response; and
(iii) \emph{Federated Operational Sovereignty}, the interoperability, supply-chain, and federation substrate that enables a unified European trust fabric.

Each cluster is analysed through a three-phase transition matrix: Phase 1, Near-Term Implementation and Pilot Actions (2026--2030), deploying near-ready solutions, enforcing enacted regulation, and closing governance gaps; Phase 2, Mid-Term Scale-Up, Industrial R\&D and Standardisation (2030--2035), industrialising sovereign capabilities, federating architectures, and building workforce and institutional infrastructure; and Phase 3, Long-Term Strategic Autonomy and Leadership (2035--2040), achieving full operational maturity, pioneering frontier research, and resolving structural doctrine questions.

For each cluster and phase, the assessment considers four perspectives: R\&I maturity; policy, regulation, and standards; operational deployment and adoption; and market uptake and strategic impact. The roadmap uses \emph{transition conditions} rather than hard gates: evidence-informed progression markers indicating when a cluster is sufficiently prepared to advance. A transition condition is a short statement that a technical basis has matured for the next stage; a policy or standardisation basis has emerged; a deployment pathway has been validated; a strategic ecosystem condition has materialised; or a Sovereignty Level improvement has become plausible through reduced dependency or increased European capacity. Transition conditions are not deterministic predictions or precise quantitative thresholds unless evidence supports them.


\section{System Model}
\label{sec:model}

The continuum is modelled as a layered reference architecture spanning
four execution environments (HPC, Cloud, Edge, Devices) and organised
into six Security Layers (vertical stack) and four Transversal
Cross-Layers (horizontal enablers). Figure~\ref{fig:framework} shows the
schema. The decomposition aligns with ETSI MEC GS~003 \cite{etsimec003},
NIST SP~800-207 and SP~800-207A \cite{rose2020zta,chandramouli2023zta},
NIST CSF 2.0 \cite{nistcsf2}, ISO/IEC 27400 \cite{iso27400}, ENISA's
baseline for IoT \cite{enisa2017iot}, and the distributed continuum
literature \cite{dustdar2023,bittencourt2018,ecroadmap2024}. Layers are
analytically distinct but operationally interdependent: a gap at one
layer propagates to adjacent layers \cite{alwarafy2021,ranaweera2021}.

\begin{figure}[t]
\centering
\fbox{\parbox{0.97\columnwidth}{\centering
\scriptsize
\setlength{\tabcolsep}{3pt}
\renewcommand{\arraystretch}{1.15}
\textit{Execution environments}\\[2pt]
\begin{tabular}{|>{\centering\arraybackslash}p{0.22\columnwidth}|>{\centering\arraybackslash}p{0.22\columnwidth}|>{\centering\arraybackslash}p{0.22\columnwidth}|>{\centering\arraybackslash}p{0.22\columnwidth}|}
\hline
\cellcolor{execbg}HPC & \cellcolor{execbg}Cloud & \cellcolor{execbg}Edge & \cellcolor{execbg}Devices \\
\hline
\multicolumn{4}{|c|}{\cellcolor{layerbg}L6: Network and Perimeter} \\
\hline
\multicolumn{4}{|c|}{\cellcolor{layerbg}L5: Data-Centric (PQC, TEE, QKD, SMPC)} \\
\hline
\multicolumn{4}{|c|}{\cellcolor{layerbg}L4: User and Orchestration (ZT, IAM, PAM)} \\
\hline
\multicolumn{4}{|c|}{\cellcolor{layerbg}L3: Application (Workloads, APIs)} \\
\hline
\multicolumn{4}{|c|}{\cellcolor{layerbg}L2: Software (OS, Middleware, OTA)} \\
\hline
\multicolumn{4}{|c|}{\cellcolor{layerbg}L1: Hardware (RoT, PhySec)} \\
\hline
\multicolumn{4}{|c|}{\cellcolor{crossbg}CL1: Cyber-Resilient SW/AI Lifecycle} \\
\hline
\multicolumn{4}{|c|}{\cellcolor{crossbg}CL2: Dynamic Risk and SecOps} \\
\hline
\multicolumn{4}{|c|}{\cellcolor{crossbg}CL3: Interoperability (Simpl, CAMARA, MIMs)} \\
\hline
\multicolumn{4}{|c|}{\cellcolor{crossbg}CL4: Compliance-by-Design and Assurance} \\
\hline
\end{tabular}
}}
\caption{Security framework: four execution environments, six Security
Layers (vertical stack, L1--L6), four Transversal Cross-Layers
(CL1--CL4).}
\label{fig:framework}
\end{figure}

\subsection{Security Layers (L1--L6)}

\textbf{L1 Hardware-Centric Security.} Physical and silicon-level trust
anchors: roots of trust, secure boot, remote attestation, TPMs, PUFs,
and physical-layer detection primitives \cite{tpm2spec,sailer2004}.
Without a verifiable silicon anchor, higher-layer attestation collapses
into unfounded assumptions. L1 is the principal locus of hardware
sovereignty and the layer where EUCC ``High'' certification
\cite{commreg2024eucc} operates as the regulatory instrument of trust.

\textbf{L2 Software-Centric Security.} Hypervisors, container runtimes,
distributed operating systems, and the update pipelines that maintain
them. Central concerns are attack-surface reduction through hardening
and cryptographically verified, failure-resilient software distribution
to field-deployed nodes where physical remediation is infeasible
\cite{ietfrfc9019,uptanestd}.

\textbf{L3 Application-Centric Security.} Containerised applications,
serverless functions, service meshes, and the APIs that expose them.
NIST SP~800-204, 204A, 204C
\cite{chandramouli2019,chandramouli2020,chandramouli2022devsecops} and
the OWASP API Security Top 10 \cite{owaspapi2023} provide the canonical
reference controls; AI Act and GDPR Article~25 obligations surface here
as runtime behaviours.

\textbf{L4 User and Orchestration-Centric Security.} Continuous,
context-adaptive authorisation aligned with NIST SP~800-207 and 207A
\cite{rose2020zta,chandramouli2023zta}; federated and self-sovereign
identity based on W3C Verifiable Credentials \cite{w3cvc2025} and, in
Europe, eIDAS~2 \cite{eueidas2} with the EU Digital Identity Wallet
implementing acts \cite{eudiwallet2024}; privileged-access controls
that eliminate standing credentials.

\textbf{L5 Data-Centric Security.} Protection of data at rest, in
transit, and in use. Confidential computing through TEEs \cite{ccc2022};
PQC based on FIPS~203/204/205
\cite{nistfips203,nistfips204,nistfips205}; signalling-data protection;
distributed trust; Quantum Key Distribution for high-value links
\cite{etsiqkd014,etsiqkd004}; Secure Multi-Party Computation for
collaborative processing across distrusting parties.

\textbf{L6 Network Infrastructure and Perimeter-Centric Security.}
5G/6G cores, disaggregated RAN, IT/OT gateways, far-edge compute, and
non-terrestrial segments \cite{tgpp33501,oran2024,snsju2023,tgppntn}.
Security integrates with physical resilience: digital breaches must not
cascade into kinetic damage.

\subsection{Transversal Cross-Layers (CL1--CL4)}

\textbf{CL1 Automated Security Processes and Cyber-Resilient AI
Lifecycle.} Secure-by-design development, runtime supply chain
verification, and continuous adversarial validation, anchored in NIST
SP~800-204C DevSecOps patterns \cite{chandramouli2022devsecops}, the NIST
AI Risk Management Framework and its Generative AI Profile
\cite{nistairmf,nistaigenai}, the SLSA provenance framework
\cite{slsa2023}, and the adversarial robustness literature
\cite{madry2018}.

\textbf{CL2 Dynamic Risk Management and Operational Cyber Resilience.}
Federated SOC operation \cite{concordia2022}, NWDAF-integrated analytics
\cite{tgpp23288}, LLM-based SecOps, and CTI lifecycle governance using
MITRE ATT\&CK \cite{mitre2020}, STIX and TAXII~2.1
\cite{stix21,taxii21}, and ENISA threat landscape reporting
\cite{enisaetl2025}.

\textbf{CL3 Interoperable Security Mechanisms.} Federation anchored in
Simpl \cite{ecsimpl}, CAMARA Network APIs \cite{camara2025}, semantic
interoperability through OASC MIMs \cite{oascmim5} and ETSI NGSI-LD
\cite{etsingsild}, and federated cyber ranges.

\textbf{CL4 Compliance-by-Design, Agile Governance, and Automated
Assurance.} Continuous certification (EUCC; EUCS when adopted)
\cite{commreg2024eucc,eucsa} using NIST OSCAL
\cite{oscal2024,iorga2024oscal}; multicloud key management; privacy-preserving
federated computation under GDPR Article~25; cyber risk quantification;
and AI trustworthiness under AI Act Article~9 \cite{euaiact,nistairmf}.


\section{Focus Areas}
\label{sec:focus}

Table~\ref{tab:focus} consolidates the focus areas per layer. Each is
the atomic unit from which the capability profiles of
Section~\ref{sec:clusters} are derived. The table replaces an
exhaustive per-layer narrative; layer-level context is provided in
Section~\ref{sec:model}.

\begin{table*}[t]
\caption{Focus areas per Security Layer and Cross-Layer.}
\label{tab:focus}
\centering
\footnotesize
\begin{tabularx}{\textwidth}{@{}lXl@{}}
\toprule
\textbf{Code} & \textbf{Focus area} & \textbf{Primary references} \\
\midrule
L1.1 & Hardware Roots of Trust; EU open-hardware (RISC-V) to EUCC ``High''
  & \cite{tpm2spec,sailer2004,opentitan2024,keystone2020,eurohpcdare2025} \\
L1.2 & Physical-layer threat detection and RF fingerprinting
  & \cite{sailer2004} \\
L2.1 & OS and middleware hardening; sovereign edge Linux
  & \cite{chandramouli2022devsecops} \\
L2.2 & Secure OTA patch and update management
  & \cite{ietfrfc9019,uptanestd} \\
L3.1 & Workload security (service mesh, mTLS, sidecar)
  & \cite{chandramouli2019,chandramouli2020} \\
L3.2 & In-band network telemetry and deep observability
  & \cite{chandramouli2020} \\
L3.3 & API security gateway and policy enforcement
  & \cite{owaspapi2023,chandramouli2022devsecops} \\
L4.1 & Zero-Trust orchestration (PEP/PDP)
  & \cite{rose2020zta,chandramouli2023zta} \\
L4.2 & Federated IAM and SSI
  & \cite{w3cvc2025,eueidas2,eudiwallet2024} \\
L4.3 & Privileged Access Management and JIT
  & \cite{rose2020zta} \\
L5.1 & Confidential computing and TEEs
  & \cite{ccc2022} \\
L5.2 & Post-Quantum Cryptography for edge/IoT
  & \cite{nistfips203,nistfips204,nistfips205,etsiqshybrid,niscg2025pqc} \\
L5.3 & Signalling/data-plane cryptographic authentication & \\
L5.4 & Distributed reputation and trust networking (DLT) & \\
L5.5 & QKD for critical links
  & \cite{etsiqkd014,etsiqkd004} \\
L5.6 & Secure Multi-Party Computation & \\
L6.1 & Telco-cloud and O-RAN security
  & \cite{tgpp33501,oran2024,oranzt2024,snsju2023} \\
L6.2 & Physical resilience and IT/OT safety & \\
L6.3 & Edge-native autonomous guardians & \\
L6.4 & Non-terrestrial network and satellite security
  & \cite{tgppntn} \\
\midrule
CL1.1 & Continuous DevSecOps
  & \cite{chandramouli2022devsecops} \\
CL1.2 & Model integrity and DataBOMs
  & \cite{nistairmf,nistaigenai} \\
CL1.3 & Dynamic supply chain verification; Runtime SBOMs
  & \cite{slsa2023} \\
CL1.4 & Automated adversarial testing and purple teaming
  & \cite{madry2018} \\
CL2.1 & Automated SOC and collaborative CTI
  & \cite{concordia2022} \\
CL2.2 & High-precision telemetry and security signalling & \\
CL2.3 & Internet of Behaviour and evidence sharing & \\
CL2.4 & AI-native analytics (NWDAF, LLMSecOps)
  & \cite{tgpp23288} \\
CL2.5 & CTI lifecycle management
  & \cite{mitre2020,stix21,taxii21,enisaetl2025} \\
CL2.6 & Deception technology and active defence & \\
CL3.1 & Simpl integration and heterogeneous interoperability
  & \cite{ecsimpl} \\
CL3.2 & Neutral host and multi-tenant isolation & \\
CL3.3 & Pan-European NaaS federation
  & \cite{camara2025} \\
CL3.4 & Semantic security interoperability
  & \cite{oascmim5,etsingsild} \\
CL3.5 & Cyber range and digital twin federation & \\
CL4.1 & PQC readiness and Zero-Trust architectures
  & \cite{rose2020zta,chandramouli2023zta} \\
CL4.2 & Continuous automated certification (EUCC; EUCS)
  & \cite{commreg2024eucc,eucsa,oscal2024,iorga2024oscal} \\
CL4.3 & Multicloud Key Management as a Service & \\
CL4.4 & Privacy-preserving federated computation
  & \cite{eugdpr} \\
CL4.6 & Cyber risk quantification and board-level reporting & \\
CL4.7 & AI trustworthiness and algorithmic accountability
  & \cite{euaiact,nistairmf} \\
\bottomrule
\end{tabularx}
\end{table*}

\section{Priority Clusters and Transition Matrices}
\label{sec:clusters}

The ten capabilities are consolidated into three priority clusters,
each presented through the seven-dimension schema with per-priority
sub-entries, followed by a three-phase transition matrix.

\subsection{Cluster~1: Sovereign Trusted Infrastructure}
\label{sec:c1}

Cluster~1 unifies P1 (Hardware Roots of Trust), P2 (PQC Agility), P3
(Lightweight Cryptography and Physical-Layer Security), P8 (Confidential
Computing), and P9 (Resilient-by-Design Infrastructure and Island Mode).
Federated policies and autonomous defence gain operational credibility
in proportion to the maturity of the execution substrate. The substrate
must guarantee hardware integrity through open-source silicon roots of
trust, quantum-resistant confidentiality through crypto-agile schemes
aligned with FIPS~203/204/205, and autonomous resilience under network
or orchestration-plane compromise. The shared critical dependency is
European sovereign fabrication capacity under the EU Chips Act
\cite{euchipsact,eurochipsaudit2025}.

\paragraph{D1 Baseline and Target.}
\emph{P1:} Infineon, NXP, and STMicroelectronics achieved EUCC ``High''
for proprietary secure element families in 2025 \cite{enisaeucc2025};
fabrication remains non-EU. No EU-origin RISC-V processor has received
EUCC or Common Criteria certification. Target: $\geq 60\%$ of critical
infrastructure edge nodes use EU-designed, EUCC ``High'' certified
secure elements by 2040.
\emph{P2:} Edge and IoT environments depend on classical algorithms.
The NIS~CG roadmap \cite{niscg2025pqc} mandates national roadmaps by
December~2026, high-risk critical infrastructure transitioned by
December~2030, full completion by December~2035. Target: quantum-native
continuum with full cryptographic agility.
\emph{P3:} Extreme-edge devices (mMTC, body-area, backscatter) are
constrained in compute, energy, and memory; conventional cryptography
is impractical. Target: lightweight, adaptive security mechanisms,
including zero-energy physical-layer authentication.
\emph{P8:} Intel TDX, AMD SEV-SNP, and ARM TrustZone dominate
operational TEE deployments; EU alternatives (Keystone
\cite{keystone2020}, OpenTitan \cite{opentitan2024}) are at research or
prototype stage. Target: end-to-end privacy-preserving AI pipelines with
progressive migration to EU-origin hardware.
\emph{P9:} Current infrastructures assume persistent connectivity.
Target: validated $\geq 72$~h Island Mode across NIS2-essential and
DORA-critical deployments.

\paragraph{D2 Primary Drivers.}
\emph{Threat:} state-sponsored exploitation of debug interfaces and
side-channel attacks on field-deployed nodes; supply chain insertion;
``Harvest~Now, Decrypt~Later'' campaigns; cascading failures from
cyberattacks, extreme weather, and geopolitical tensions.
\emph{Regulatory:} CRA Articles~10--11 with enforcement from
December~2027 \cite{eucra2024}; NIS2 Article~21(2)(e) on cryptographic
policies \cite{eunis2}; AI Act Articles~9 and~15 on hardware integrity
for high-risk systems \cite{euaiact}; NIS~CG PQC roadmap
\cite{niscg2025pqc}; DORA Article~11 on business continuity
\cite{eudora}; SNS~JU SRIA on 6G resilience \cite{snsju2023}.

\paragraph{D3 High-Impact Activities.}
Phase~1: fund three to five RISC-V secure element pilots under Horizon
Europe Cluster~3 targeting EUCC ``High'' pre-evaluation; submit
OpenTitan and Keystone prototypes; optimise FIPS~203/204/205 for
8/16-bit microcontrollers; deploy hybrid ML-KEM plus ECDH-P384 on the
cloud side per ETSI TS~103~744 \cite{etsiqshybrid}; validate Island
Mode soft-state protocols in 5G testbeds. Phase~2: first EUCC ``High''
certified EU-manufactured RISC-V secure elements in production via
Chips~Act pilot lines (NanoIC, FAMES); RISC-V Keystone TEEs in
multi-provider Data Spaces; certified 72~h Island Mode under 3GPP/ETSI
fallback profiles in cross-border 5G corridors; mandatory hybrid PQC in
DORA-scoped deployments. Phase~3: EU-designed RISC-V secure elements as
the hardware default; FHE end-to-end from IoT sensor to HPC compute
node; zero-energy PHY authentication for battery-less mMTC; mandatory
validated 72~h autonomous Island Mode for NIS2-essential and
DORA-critical deployments.

\paragraph{D4 Catalysts and Enablers.}
EU Chips Act pilot lines (NanoIC, FAMES, APECS, PIXEurope, WBG;
combined funding $\approx \mathrm{\text{\euro}}3.7$B \cite{euchipsact});
DARE~SGA1 (EuroHPC, $\mathrm{\text{\euro}}240$M) \cite{eurohpcdare2025};
ECCC Strategic Agenda \cite{eccc2024}; CEN/CENELEC JTC~13 hardware
security standards; NIST FIPS~203/204/205
\cite{nistfips203,nistfips204,nistfips205} and ETSI TS~103~744
\cite{etsiqshybrid}; SNS~JU 6G resilience projects \cite{snsju2023}.

\paragraph{D5 Disruptors and Systemic Risks.}
EU Chips Act fabrication capacity shortfall: ECA Special Report 12/2025
projects $\approx 11.7\%$ global semiconductor share by 2030 against the
20\% target \cite{eurochipsaudit2025}; the cancellation of Intel
Magdeburg in July~2025 removed the largest EU-based advanced node
commitment. Early arrival of a Cryptographically Relevant Quantum
Computer before migration is complete. PQC computational overhead on
battery-powered IoT (ML-KEM at $\approx 1.5\times$ ECDH on
Cortex-M4 at equivalent security \cite{nistir8413}). Side-channel
breakthroughs against current TEE architectures. Insufficient
cryptographic inventory of legacy OT/IoT systems.

\paragraph{D6 Failsafe Mechanisms.}
TEE abstraction layers to avoid vendor lock-in during the Phase~1
dependency period; hybrid cryptographic schemes (classical plus PQC) for
backward compatibility and harvest resistance; crypto-agility frameworks
enabling sub-second algorithm switching; diversification across PQC
algorithm families (lattice, code, hash) to avoid single-point
cryptanalytic failure; autonomous node quarantine via eBPF network
policy if attestation fails; layered lightweight cryptography with PHY
signatures as last-resort authentication under energy starvation;
redundant trust anchors, segmented autonomy zones, and tamper-resistant
local decision logic for Island Mode.

\paragraph{D7 Transversal Meta-Layer.}
\emph{Sovereignty:} eliminates proprietary hardware dependencies and
quantum decryption exposure; current TEE dependency approaches 100\%.
\emph{Trust and Ethics:} EUCC ``High'' aligns with Common Criteria
EAL4+, AVA\_VAN.5; PQC alignment with GDPR Article~32
\cite{eugdpr}; the 18-Member-State joint statement ``Securing
Tomorrow, Today'' (November~2024) formalises PQC as a collective
sovereignty concern.
\emph{Sustainable Security:} energy-efficient secure silicon for
far-edge; lightweight PQC extends device lifecycles and avoids premature
replacement.

\begin{table*}[t]
\caption{Cluster~1 transition matrix, Sovereign Trusted Infrastructure.}
\label{tab:c1}
\centering
\footnotesize
\begin{tabularx}{\textwidth}{@{}lXXX@{}}
\toprule
 & \textbf{Phase~1 (2026--2030)} & \textbf{Phase~2 (2030--2035)} & \textbf{Phase~3 (2035--2040)} \\
\midrule
\textbf{ODP} & Applied R\&D to Fragmented Piloting & Standardised Integration & Ubiquitous Operation \\
\textbf{SSL} & High Foreign Dependency & Emerging EU Ecosystem & Full Sovereign Autonomy \\
\textbf{Operational State} &
RISC-V secure element pilots in three to five Horizon Europe consortia;
Intel/AMD/ARM TEEs dominate all operational deployments; classical
cryptography on far-edge IoT; Island Mode validated in 5G testbeds by
SNS~JU consortia. &
First EUCC ``High'' EU-manufactured RISC-V secure elements in production
from Chips~JU pilot lines; deployed in energy substations and telecom
critical infrastructure in $\geq 5$ Member States; RISC-V Keystone TEEs
operational in multi-provider Data Spaces alongside Intel TDX; 72~h
Island Mode certified in cross-border 5G corridors. &
EU-designed secure elements as default in new EU ICT products;
$\geq 60\%$ of critical infrastructure edge nodes on EUCC ``High''
EU-designed elements; FHE end-to-end across continuum;
72~h autonomous Island Mode mandatory in NIS2-essential and
DORA-critical deployments. \\
\textbf{Critical Enabler} &
EUCC ``High'' certification of an EU-origin open-source RISC-V secure
element before end of phase. &
Ratification of ETSI TS profiles for FIPS~203/204/205 tailored to EU
deployment contexts; operational Chips~JU fabrication capacity at the
advanced nodes required for secure elements. &
Mandatory EUCC ``High'' without waiver for new EU ICT products; PQC-native
silicon as EU procurement default. \\
\bottomrule
\end{tabularx}
\end{table*}

\subsection{Cluster~2: Cognitive Security and Autonomous Cyber Defence}
\label{sec:c2}

Cluster~2 encompasses P6 (LLMSecOps) and P10 (Proactive APT Management),
with cross-layer dependencies on Model Integrity and DataBOMs (CL1.2),
Dynamic Supply Chain Verification (CL1.3), and AI Trustworthiness
(CL4.7). The cluster is structured by three inseparable sub-dimensions:
AI as a security instrument (sovereign models, edge-native agents,
autonomous SOC, Moving Target Defence); AI as adversary amplifier
(polymorphic malware, automated exploit discovery); AI as target
(training pipelines and inference endpoints as attack surfaces under
AI Act Article~9).

\paragraph{D1 Baseline and Target.}
\emph{P6:} US-platform SIEMs dominate the EU market, no
EU-headquartered vendor in the Gartner Leaders quadrant
\cite{gartner2025siem}. Sovereign EU security-specific foundation models
are absent; OpenEuroLLM \cite{openeurollm2025} has no dedicated security
stream. Target: cognitive SecOps model with EU-sovereign models,
auditable under AI Act Article~9, covering $\geq 50\%$ of new
critical infrastructure SOC deployments by 2040.
\emph{P10:} Average APT dwell times in European organisations are
150--200 days. Target: $\leq 72$~h dwell time in Phase~2 and
$\leq 24$~h in Phase~3 across NIS2-essential deployments.

\paragraph{D2 Primary Drivers.}
\emph{Threat:} weaponisation of generative AI for polymorphic malware,
automated exploit discovery, and scaled social engineering;
nation-state actors blending espionage with organised crime;
memory-only execution and abuse of legitimate administrative tools.
\emph{Regulatory:} AI Act Articles~9, 13, and 15 creating mandatory
requirements for autonomous security AI classified as high-risk
\cite{euaiact}; NIS2 Article~21(2)(b) on incident handling
\cite{eunis2}. Conformity assessment capacity lags obligations by
approximately 24--36 months \cite{ep2026aiact}.

\paragraph{D3 High-Impact Activities.}
Phase~1: integrate commercial LLMs into SOC workflows for alert triage
and MITRE ATT\&CK mapping; route NWDAF telemetry into centralised AI
analytics in three to five 5G SA pilots \cite{tgpp23288}; train a first
EU-sovereign SLM ($\leq 10$B parameters) on GDPR-compliant security
corpora using EuroHPC \cite{eurohpcaifactories}; apply ML-based UEBA for
APT IoC detection; target sub-72~h dwell time in one certified
NIS2-essential pilot. Phase~2: deploy EU-sovereign SLMs at the edge for
real-time SOC analytics; deception infrastructure (honeynets, credential
canaries) for calibration; autonomous triage under ETSI GS~ZSM~009
\cite{etsizsm}; Zero-Trust reconfigurations per NIST~SP~800-207
\cite{rose2020zta}. Phase~3: swarm-based agentic AI for decentralised
threat intelligence; continuous Moving Target Defence (IP, topology,
software diversity); continuous automated adversarial robustness
validation.

\paragraph{D4 Catalysts and Enablers.}
EuroHPC AI Factories (19 confirmed across 16 Member States; cumulative
funding $>\mathrm{\text{\euro}}2.7$B; IT4LIA, PIAST, MeluXina-AI list
cybersecurity as a vertical) \cite{eurohpcaifactories}; ETSI GS~ZSM~009
\cite{etsizsm}; ECCC Strategic Agenda \cite{eccc2024}; MITRE ATT\&CK
\cite{mitre2020} and STIX/TAXII~2.1 \cite{stix21,taxii21}; ENISA threat
landscape reporting \cite{enisaetl2025}; NIST AI RMF and Generative AI
Profile \cite{nistairmf,nistaigenai}.

\paragraph{D5 Disruptors and Systemic Risks.}
Hallucinations and automation bias in safety-relevant contexts;
adversarial prompt manipulation and data leakage into models; absence of
a standardised adversarial robustness test suite for security AI (the
single most consequential standards gap in the cluster); intelligence
fragmentation and weak information-sharing incentives; adversarial
adaptation outpacing defensive innovation; CA body capacity shortfall
\cite{ep2026aiact}.

\paragraph{D6 Failsafe Mechanisms.}
Human-in-the-loop validation for Tier-1 critical actions; bounded
autonomy governed by ETSI GS~ZSM~009 \cite{etsizsm}; shadow-mode
rollback if autonomous agents exhibit unstable control loops;
policy-constrained orchestration with audit logging; autonomous
real-time micro-segmentation to quarantine compromised components
under Zero-Trust; layered detection, deception, and rapid recovery
playbooks.

\paragraph{D7 Transversal Meta-Layer.}
\emph{Sovereignty:} reduces approximately 80\% US-platform dependency
in enterprise SIEM.
\emph{Trust and Ethics:} AI Act Article~9 compliance-by-design,
explainability, and accountability for autonomous security decisions;
NIS2 sector-specific implementing acts should reference MITRE ATT\&CK
\cite{mitre2020}.
\emph{Sustainable Security:} edge-native SLMs reduce inference-time
energy cost versus cloud round-trips for high-volume security
telemetry.

\begin{table*}[t]
\caption{Cluster~2 transition matrix, Cognitive Security and Autonomous
Cyber Defence.}
\label{tab:c2}
\centering
\footnotesize
\begin{tabularx}{\textwidth}{@{}lXXX@{}}
\toprule
 & \textbf{Phase~1 (2026--2030)} & \textbf{Phase~2 (2030--2035)} & \textbf{Phase~3 (2035--2040)} \\
\midrule
\textbf{ODP} & Applied R\&D, Siloed AI Copilots & Standardised Integration & Ubiquitous Operation \\
\textbf{SSL} & High Reliance on US Foundation Models & Sovereign EU Edge LLMs Emerging & Global Leadership in Trustworthy AI Security \\
\textbf{Operational State} &
Commercial LLMs integrated into SOC workflows for triage; NWDAF pilots
in three to five 5G SA deployments; first EU-sovereign SLM under
training on EuroHPC AI Factory compute; sub-72~h dwell time achieved in
at least one certified pilot; no operational autonomous containment
agent. &
EU-sovereign SLMs deployed at the edge for real-time SOC analytics
without cloud inference; autonomous triage governed by ETSI GS~ZSM~009
in production; sub-72~h dwell time in certified critical infrastructure
deployments; deception infrastructure operational for calibration. &
Swarm-based agentic AI defending the continuum via self-evolving
decentralised CTI; Moving Target Defence continuous in certified
environments; $\leq 24$~h dwell time across NIS2-essential
deployments; EU-sovereign AI security tools in $\geq 50\%$ of new
critical infrastructure SOC deployments. \\
\textbf{Critical Enabler} &
CEN/CENELEC or ETSI TC CYBER publication of a standardised adversarial
robustness test suite for security AI before end of phase. &
Operationalisation of an AI Act Article~9 conformity assessment
framework specific to autonomous security agents; activation of EuroHPC
AI Factory compute contracts for EU-sovereign security model training. &
International adoption of EU AI Act frameworks for trustworthy
autonomous security agents by at least one non-EU partner jurisdiction. \\
\bottomrule
\end{tabularx}
\end{table*}

\subsection{Cluster~3: Federated Operational Sovereignty}
\label{sec:c3}

Cluster~3 encompasses P4 (Security Interoperability via Simpl), P5
(Supply Chain Verification), and P7 (Pan-European NaaS Federation). The
cluster addresses three mechanisms that transform sovereign components
into a collectively sovereign continuum: federated interoperability
through Minimal Interoperability Mechanisms (MIMs) in the Simpl stack
\cite{ecsimpl} under Data~Act portability mandates \cite{eudataact};
supply chain integrity through Runtime SBOMs, Digital Product
Passports, and cryptographic proof-of-execution under CRA Annex~I
\cite{eucra2024}; and compute federation through harmonised CAMARA
Network APIs \cite{camara2025} and cross-border edge roaming. Supply
chain verification bridges Clusters~1 and~3 because hardware attestation
depends on the RISC-V secure elements matured in Cluster~1.

\paragraph{D1 Baseline and Target.}
\emph{P4:} European cloud-edge ecosystems remain fragmented across
proprietary vendor silos with incompatible security policies. The
Simpl-Open iterative MVP has been available since January~2025, with
enhanced releases in June and December~2025 \cite{ecsimpl}. OASC MIM6
(Security management) remains a ``Work item'' without a formally
published specification. Target: unified European trust fabric based on
Simpl MIMs ratified as CEN/CENELEC standards and referenced in the EU
Cloud Rulebook.
\emph{P5:} Current supply chains are opaque and globally distributed.
CRA Annex~I Part~II mandates static SBOM disclosure from December~2027
\cite{eucra2024}; NIS2 Article~21(2)(h) mandates supply chain risk
assessment \cite{eunis2}; DORA Article~28 requires continuous ICT
third-party monitoring \cite{eudora}, operationalised through the first
19 Critical ICT Third-Party Provider designations on 18~November~2025
\cite{esa2025ctpp}. Target: progression from static SBOM disclosure
(Phase~1) to Runtime SBOMs with cryptographic proof-of-execution
(Phase~2) to ZKP-verified supply chain attestation at transaction scale
(Phase~3).
\emph{P7:} The CAMARA Fall~2025 meta-release comprises 60 APIs (10
stable) \cite{camara2025}; GSMA Open Gateway reaches 86 operator groups
\cite{gsmaopengateway2026}. Cross-border edge compute roaming does not
exist at operational scale. Target: commercial EU NaaS federation as a
unified jurisdiction for data workloads, with non-EU hyperscaler market
share in EU regulated cloud procurement below 40\% for new deployments
by 2040.

\paragraph{D2 Primary Drivers.}
\emph{Threat:} SolarWinds-class software supply chain attacks;
disaggregation of telco clouds introducing multi-vendor complexity; the
sovereignty cliff where absence of a Digital Single Market prevents EU
operators from scaling against US and China hyperscalers.
\emph{Regulatory:} Data Act (applicable from 12~September~2025)
\cite{eudataact}; NIS2 coordinated incident response
\cite{eunis2}; CRA SBOM obligations (December~2027)
\cite{eucra2024}; DORA CTPP monitoring \cite{eudora}; the Digital
Networks Act proposal (COM(2026)~16, January~2026) signals EC intent
for network federation but does not mandate API standards.

\paragraph{D3 High-Impact Activities.}
Phase~1: Simpl security MIM pilots in three Common European Data Spaces
(Manufacturing, Health \cite{ehds2025}, Energy); deploy harmonised
CAMARA APIs in two
to three bilateral 5G corridor agreements; submit CRA-mandated static
SBOM declarations; initiate Digital Product Passport pilots; advance
ZKP verification towards demonstrator readiness. Phase~2: automated
policy translation via standardised Simpl MIMs across active Data
Spaces; Runtime SBOMs with cryptographic proof-of-execution in critical
infrastructure in $\geq 5$ Member States; automatic quarantine of
unverified components in isolated network slices; CAMARA cross-border
roaming scaled to $\geq 5$ operator pairs; ``Data Visa'' mechanisms for
dynamic AI model migration. Phase~3: Simpl-based security as universal
Data Space standard; ZKP-verified supply chain attestation at real-time
scale; operational EU NaaS federation; ``Schengen for Cloud'' legal
framework resolving US CLOUD Act jurisdictional conflict.

\paragraph{D4 Catalysts and Enablers.}
Digital Europe Programme funding for Simpl \cite{ecsimpl}; CEN-CLC/JTC
25 for data-space standards; Eclipse Dataspace Protocol;
IPCEI-CIS ICRA Reference Architecture v2.0 (January~2026, total
$\mathrm{\text{\euro}}2.6$B, 19 participants, 12 Member States)
\cite{ipceicis2023}; EU Ecodesign Regulation Digital Product Passport;
DORA Delegated Regulation (EU) 2024/1502 for CTPP assessment
\cite{eudora}; SNS~JU \cite{snsju2023}; Aduna as primary aggregator.

\paragraph{D5 Disruptors and Systemic Risks.}
Vendor resistance to open standards; weakest-link failures in federated
systems; uneven Member State compliance maturity; US CLOUD Act
jurisdictional conflict \cite{microsoftsenat2025}; concentration of
manufacturing dependencies; limited multi-tier auditability; SME
exclusion from compliance pathways; hyperscaler platform dominance
through proprietary API ecosystems (approximately 70\% EU cloud share)
\cite{synergy2025}.

\paragraph{D6 Failsafe Mechanisms.}
Autonomous unfederation or isolation of nodes failing the shared
baseline; Zero-Trust architectures where federation does not imply
implicit trust \cite{rose2020zta}; cryptographic attestation of policy
enforcement; graceful isolation of unverified components into
quarantined network slices; multi-vendor diversification;
trusted fallback components; graceful degradation redirecting workloads
to national infrastructure; regional autonomy zones; multi-operator
fallback routing.

\paragraph{D7 Transversal Meta-Layer.}
\emph{Sovereignty:} breaks hyperscaler lock-in and resolves structural
CLOUD~Act conflict.
\emph{Trust and Ethics:} transparent and verifiable federation
mechanisms enabling cross-border trust; the 2025 DINUM Bilan reports
70\% European-provider share in French public cloud procurement
market-wide and 99\% within the State perimeter \cite{dinum2026bilan}.
\emph{Sustainable Security:} shared infrastructure reduces duplication
of compute and certification overhead across Member States.

\begin{table*}[t]
\caption{Cluster~3 transition matrix, Federated Operational Sovereignty.}
\label{tab:c3}
\centering
\footnotesize
\begin{tabularx}{\textwidth}{@{}lXXX@{}}
\toprule
 & \textbf{Phase~1 (2026--2030)} & \textbf{Phase~2 (2030--2035)} & \textbf{Phase~3 (2035--2040)} \\
\midrule
\textbf{ODP} & Fragmented Piloting & Standardised Integration & Ubiquitous Operation \\
\textbf{SSL} & Mixed: Emerging EU Ecosystem / Hyperscaler Dominance & Emerging EU Ecosystem & Full Federated Sovereign Autonomy \\
\textbf{Operational State} &
Simpl security MIM pilots in three Common European Data Spaces with
manual policy coordination; CAMARA APIs in two to three bilateral 5G
corridors; no operational cross-border edge compute roaming; CRA-mandated
static SBOMs submitted as point-in-time artefacts; DPP pilots initiated;
hyperscalers retain approximately 70\% of EU cloud market. &
Automated policy translation via standardised Simpl MIMs across active
Data Spaces; Runtime SBOMs with cryptographic proof-of-execution in
critical infrastructure in $\geq 5$ Member States; ZKP batch
verification operational; CAMARA roaming across $\geq 5$ operator
pairs; ``Data Visa'' mechanisms operational; non-hyperscaler share
$\geq 30\%$ in new regulated-sector procurement. &
Simpl-based security as universal cross-provider Data Space standard;
ZKP-verified supply chain attestation at real-time scale; operational EU
NaaS federation; ``Schengen for Cloud'' framework enacted; non-EU
hyperscaler share below 40\% in EU regulated cloud procurement for new
deployments; EU supply chain standards adopted in ISO/IEC and by at
least two partner jurisdictions. \\
\textbf{Critical Enabler} &
Simpl-Open iterative MVP consolidation with OASC MIM6 (Security
management) progressing from Work item to formally published
specification. &
Ratification of Simpl security MIMs as CEN/CENELEC standards referenced
as mandatory in the EU Cloud Rulebook; completion of DORA continuous
ICT third-party monitoring operationalisation. &
Enactment of the ``Schengen for Cloud'' legal framework resolving US
CLOUD Act jurisdictional conflict. \\
\bottomrule
\end{tabularx}
\end{table*}

\subsection{Cross-Cutting Enablers}
\label{sec:crosscut}

Four dependencies span at least two clusters and require dedicated
governance to avoid becoming systemic bottlenecks.

\emph{EU Chips Act sovereign fabrication} (Clusters~1 and~3). Five
Chips~JU pilot lines are operational, with combined EU and national
funding of $\mathrm{\text{\euro}}3.7$B \cite{euchipsact}. No pilot line
is dedicated to secure element fabrication; restriction of access to
any major non-EU foundry would immediately impair EU critical
infrastructure attestation capability. Mitigation: strategic
stockpiling, second-source agreements, and dedicated pilot lines for
security-critical components.

\emph{AI Act conformity assessment infrastructure} (Cluster~2). A 24--36
month gap between AI Act obligations and CA body operational readiness
is documented \cite{ep2026aiact}. As of March~2026, 8 of 27~Member
States have designated single points of contact. Mitigation:
accelerated CEN/CENELEC adversarial robustness test suite development
and AI Act sandbox provisions for security use cases.

\emph{Runtime SBOM attestation} (Clusters~1 and~3). Static SBOMs
(CRA-mandated) are insufficient for dynamic continuum operations.
Progression to cryptographic proof-of-execution and ZKP batch
verification is a shared Phase~2--3 transition condition.

\emph{SME access and inclusion} (all clusters). Cost barriers exclude
SMEs from sovereign hardware alternatives, AI-powered SOC tooling, and
Simpl MIM compliance. Digital Europe Programme co-funding should target
shared-infrastructure pathways and subsidised certification programmes.

\section{Strategic Recommendations}
\label{sec:recs}

The 2040 destination is a secure, sovereign, interoperable, and
resilient European computing continuum. Six actions follow from the
analysis.

\emph{Sovereign secure silicon.} Combine Horizon Europe support with
the EU Chips Act to advance open-source RISC-V secure elements and
pilot fabrication lines reaching EUCC ``High'' assurance. The 2025
certifications achieved by Infineon, NXP, and STMicroelectronics
\cite{enisaeucc2025} demonstrate framework readiness; the remaining gap
is open-source, EU-controlled fabrication.

\emph{Operationalised PQC migration.} Deliver crypto-agility middleware
for edge and IoT aligned with the NIS~CG PQC roadmap milestones
\cite{niscg2025pqc} and ETSI TS~103~744 hybrid key exchange profiles
\cite{etsiqshybrid}. Make hybrid cryptographic deployment the default
for new critical infrastructure.

\emph{Security interoperability through Simpl.} Move OASC MIM6 (Security
management) from Work item status to a published specification, then to
a ratified CEN/CENELEC standard referenced in the EU Cloud Rulebook, so
that interoperability carries procurement-level enforcement.

\emph{Continuous certification.} Use EUCC \cite{commreg2024eucc} as the
foundation and unblock EUCS \cite{eucsa}, plausibly through the 2026
Cybersecurity Act revision (COM(2026)~13) \cite{ec2025cloudsov}, to
extend continuous assurance to cloud services.

\emph{Federated SOC and CTI.} Build on DORA CTPP oversight
\cite{esa2025ctpp} and cross-border CSIRT cooperation
\cite{concordia2022} to support pan-European situational awareness.

\emph{SME inclusion mandate.} Implement tiered certification,
shared cyber range access, and reusable compliance tooling so that
NIS2, CRA, and AI Act compliance does not become a structural barrier
to innovation.

\section{Discussion}
\label{sec:disc}

\subsection{Validation}
The roadmap is treated as an evolving strategic artefact benefiting
from expert review, stakeholder consultation, and iterative revision.
Validation is conducted through expert workshops within the NexusForum
and EuCloudEdgeIoT Working Group on Cybersecurity; consultation with
national cybersecurity agencies and industry; and policy dialogue with
European Commission services. Validation addresses completeness and
relevance of the architectural decomposition, clarity and non-overlap
of capability formulations, coherence of priority clustering,
credibility of transition conditions, and reasonableness of Sovereignty
Level assessments.

\subsection{Limitations}
The architectural decomposition reflects a particular analytical
perspective and may not capture intersections with adjacent domains
(energy, industrial policy, trade). The seven-dimension profiles rely
on expert judgement rather than empirical measurement. Transition
conditions sacrifice precision for robustness; they are suited to
strategic orientation rather than operational planning. The Sovereignty
Level construct resists precise quantification and is used
interpretively. The paper also reflects the ambition--implementation
gap of the 2026 regulatory landscape: NIS2 transposition delays in
several large Member States \cite{ec2025nis2status}, a contested 20\%
semiconductor target \cite{eurochipsaudit2025}, EUCS remaining
unadopted \cite{ec2025cloudsov}, and no harmonised CRA standards yet
cited in the Official Journal. These realities constrain Phase~1
governance specificity and locate the roadmap's primary value in
structuring coordination.

\subsection{Methodological Positioning}
The methodology is suited to strategic agenda setting, expert
discussion, cross-priority comparison, and iterative roadmap evolution
under conditions of ecosystem uncertainty. Its value lies in providing
a shared analytical structure within which stakeholders can coordinate
research, policy, and investment decisions toward a common set of
objectives.

\section{Conclusion}
\label{sec:concl}

The Cloud--Edge--IoT continuum requires Europe to move beyond reactive
and fragmented cybersecurity models towards autonomous, federated, and
certified Security-by-Design. The structured methodology presented here
identifies security dimensions, derives strategic capabilities, and
consolidates them into priority clusters aligned with resilience and
digital sovereignty objectives. The ten capabilities form an
interdependent strategic portfolio that must be developed in parallel.
The 2025--2026 period has produced concrete milestones (EUCC ``High''
certifications across all three major EU secure element manufacturers,
a binding PQC migration framework, DORA operationalisation through CTPP
designation, and emerging Simpl infrastructure) and persistent gaps
(NIS2 transposition delays, EUCS deadlock, contested semiconductor
targets, continued hyperscaler dominance at approximately 70\% of the
EU cloud market). The ambition--implementation gap defines where the
coordination challenge is most acute and where strategic investment
will yield the greatest returns. Future work should extend this
analysis through empirical validation of Sovereignty Level assessments,
sector-specific instantiation (energy, healthcare, financial services),
and integration with the next EU Framework Programme funding landscape.

\bibliographystyle{IEEEtran}
\bibliography{references}

\EOD
\end{document}